\section{Simulation Parameters and Benchmarks}\label{sec:sim}

\begin{table}[t]
\centering
\caption{Default Simulation Parameters}
\label{tab:params}
\begin{tabular}{@{}ll@{}}
\toprule
Parameter & Value \\
\midrule
Orbit altitude & 600 km \\
Carrier frequency & 2 GHz (S band) \\
Satellite antennas $N$ & $16 \times 16$ UPA \\
STAR-RIS elements $M$ & 256 (128 in runtime demo) \\
Rician factor $\kappa$ (sat link) & 10 dB \\
Satellite antenna gain $G_s$ & 45 dBi \\
Users ($K_t, K_r$) & (2, 2) \\
Satellite power $P_{\max}$ & 50 dBm \\
Spoofer power $P_a$ & $-10$ dBm (GNSS-class spoofing) \\
Noise power $\sigma^2$ & $-110$ dBm \\
CRB threshold $\epsilon_{\mathrm{CRB}}$ & $1\,\mathrm{m}^2$ \\
Element spacing $d$ & $\lambda/2$ \\
AN power fraction & 20\% (isotropic) \\
Spoofer--user separation & 50 m \\
Direct-link shadowing (no-RIS benchmark) & 30 dB \\
\bottomrule
\end{tabular}
\end{table}

Table~\ref{tab:params} lists the default parameters. With $M=256$ and $d=\lambda/2$, the STAR-RIS aperture is $D \approx 19$\,m and the corresponding Rayleigh distance $2D^2/\lambda \approx 4.9$\,km, so ground users within tens of meters lie deep inside the near-field region, justifying the spherical-wave model of~\eqref{eq:nfsteer}.

